\documentclass[aspectratio=169, 11pt]{beamer}

\usepackage[utf8]{inputenc}
\usepackage[T1]{fontenc}
\usepackage[brazil]{babel}
\usepackage{listings}
\usepackage{xcolor}
\usepackage{tikz}

% ---------------- PALETA DE CORES (ESTILO DJANGO) ----------------
\definecolor{djangoDark}{HTML}{092E20}    % Verde escuro institucional Django
\definecolor{djangoMid}{HTML}{0C4B33}     % Verde médio
\definecolor{djangoAccent}{HTML}{44B78B}  % Verde claro de destaque
\definecolor{bgLight}{HTML}{F8FAFC}       % Fundo claro sutil
\definecolor{textDark}{HTML}{1E293B}      % Cinza grafite para texto
\definecolor{textMuted}{HTML}{64748B}     % Cinza suave para metadados
\definecolor{codeBg}{HTML}{F1F5F9}        % Fundo dos blocos de código
\definecolor{codeBorder}{HTML}{CBD5E1}    % Borda dos blocos de código

% ---------------- CONFIGURAÇÃO DO TEMA ----------------
\usetheme{default}
\usefonttheme{professionalfonts}

\setbeamercolor{normal text}{fg=textDark, bg=white}
\setbeamercolor{structure}{fg=djangoDark}
\setbeamercolor{title}{fg=white}
\setbeamercolor{subtitle}{fg=djangoAccent}
\setbeamercolor{author}{fg=white}
\setbeamercolor{institute}{fg=codeBorder}
\setbeamercolor{date}{fg=codeBorder}

\setbeamercolor{frametitle}{fg=djangoDark, bg=white}
\setbeamercolor{block title}{fg=white, bg=djangoMid}
\setbeamercolor{block body}{fg=textDark, bg=bgLight}
\setbeamercolor{item}{fg=djangoMid}
\setbeamercolor{enumerate item}{fg=djangoMid}

% Fontes e espaçamentos
\setbeamerfont{title}{size=\Huge, series=\bfseries}
\setbeamerfont{subtitle}{size=\large}
\setbeamerfont{frametitle}{size=\Large, series=\bfseries}
\setbeamerfont{block title}{size=\normalsize, series=\bfseries}
\setbeamertemplate{navigation symbols}{}
\setbeamertemplate{itemize items}[circle]
\setbeamertemplate{blocks}[rounded][shadow=false]

% ---------------- CABEÇALHO E RODAPÉ ----------------
\setbeamertemplate{frametitle}{
  \vspace{0.35cm}
  \begin{beamercolorbox}[wd=\textwidth, leftskip=0.2cm]{frametitle}
    {\small\color{djangoAccent}\bfseries\insertsectionhead}\\[-2pt]
    \insertframetitle
  \end{beamercolorbox}
  \vspace{0.1cm}
  {\color{djangoAccent}\hrule height 1.2pt}
}

\setbeamertemplate{footline}{
  \leavevmode%
  \hbox{%
    \begin{beamercolorbox}[wd=.75\paperwidth, ht=2.6ex, dp=1.1ex, leftskip=0.6cm]{author in head/foot}%
      \color{textMuted}\scriptsize \insertshorttitle \quad $\cdot$ \quad \insertshortauthor
    \end{beamercolorbox}%
    \begin{beamercolorbox}[wd=.25\paperwidth, ht=2.6ex, dp=1.1ex, rightskip=0.6cm]{date in head/foot}%
      \hfill\color{djangoDark}\scriptsize\bfseries \insertframenumber{} / \inserttotalframenumber
    \end{beamercolorbox}%
  }%
  \vskip2pt%
}

% ---------------- SLIDE DE CAPA PERSONALIZADO ----------------
\setbeamertemplate{title page}{
  \begin{tikzpicture}[remember picture, overlay]
    \fill[djangoDark] (current page.south west) rectangle (current page.north east);
    \fill[djangoAccent] (current page.south west) rectangle ([xshift=0.35cm]current page.north west);
    \fill[djangoMid] ([yshift=1.4cm]current page.south west) rectangle (current page.south east);
  \end{tikzpicture}
  \vspace{1.2cm}
  \begin{flushleft}
    {\color{djangoAccent}\small\bfseries SEMINÁRIO}\\[0.25cm]
    {\usebeamerfont{title}\usebeamercolor[fg]{title}\inserttitle}\\[0.2cm]
    {\usebeamerfont{subtitle}\usebeamercolor[fg]{subtitle}\insertsubtitle}\\[1.2cm]
    {\color{white}\normalsize\insertauthor}\\[0.15cm]
    {\color{codeBorder}\small\insertinstitute}
  \end{flushleft}
  \vfill
  \begin{flushright}
    {\color{white}\footnotesize\insertdate}
  \end{flushright}
}

% ---------------- CONFIGURAÇÃO DE CÓDIGO (LISTINGS) ----------------
\lstset{
  backgroundcolor=\color{codeBg},
  basicstyle=\ttfamily\footnotesize\color{textDark},
  keywordstyle=\color{djangoMid}\bfseries,
  commentstyle=\color{textMuted}\itshape,
  stringstyle=\color{teal},
  breaklines=true,
  frame=leftline,
  framerule=2.5pt,
  rulecolor=\color{djangoAccent},
  xleftmargin=6pt,
  framexleftmargin=6pt,
  language=Python,
  showstringspaces=false,
  tabsize=4
}

% ---------------- METADADOS ----------------
\title{Django}
\subtitle{Framework web em Python + CRUD na prática}
\author{Emerson Gomes \texorpdfstring{$\cdot$}{-} Entony Richard \texorpdfstring{$\cdot$}{-} Guilherme Pontes \texorpdfstring{$\cdot$}{-} Matheus Henrique}
\institute{Instituto Federal de Alagoas}
\date{\today}

\begin{document}

% Capa (sem cabeçalho/rodapé)
{
\setbeamertemplate{footline}{}
\begin{frame}[plain]
  \titlepage
\end{frame}
}

% ---------------- ROTEIRO ----------------
\section{Apresentação}
\begin{frame}{Roteiro do Seminário}
  \vspace{0.2cm}
  \begin{enumerate}
    \setlength\itemsep{0.6em}
    \item \textbf{O que é o Django?} --- Conceitos, história e vantagens
    \item \textbf{Arquitetura MVT} --- Como o framework organiza a aplicação
    \item \textbf{Instalação e Estrutura} --- Criando um projeto do zero
    \item \textbf{CRUD na Prática} --- Lista de tarefas passo a passo
    \item \textbf{Demonstração e Conclusão} --- Aplicação funcionando
  \end{enumerate}
\end{frame}

% ---------------- PESSOA 1 ----------------
\section{Introdução}

\begin{frame}{O que é o Django?}
  \begin{itemize}
    \setlength\itemsep{0.55em}
    \item Framework web de alto nível escrito em \textbf{Python}.
    \item Criado em 2005 e adotado por plataformas como \textbf{Instagram}, \textbf{Pinterest} e \textbf{Spotify}.
    \item Segue a filosofia \textbf{DRY} (\textit{Don't Repeat Yourself}).
    \item Abordagem \textit{"Batteries Included"}: já traz painel administrativo, autenticação, ORM e segurança.
    \item Gratuito, maduro e de código aberto (\textit{open source}).
  \end{itemize}
\end{frame}

\begin{frame}{Por que escolher o Django?}
  \begin{columns}[T]
    \column{0.48\textwidth}
    \begin{block}{Produtividade e Escala}
      \begin{itemize}
        \item Desenvolvimento ágil com menos código repetitivo.
        \item Escala de protótipos rápidos a sistemas de alto tráfego.
        \item Sintaxe limpa e legível do Python.
      \end{itemize}
    \end{block}

    \column{0.48\textwidth}
    \begin{block}{Segurança e Ecossistema}
      \begin{itemize}
        \item Proteção nativa contra \textbf{SQL Injection}, \textbf{CSRF} e \textbf{XSS}.
        \item Documentação oficial extensa e didática.
        \item Grande comunidade e pacotes prontos.
      \end{itemize}
    \end{block}
  \end{columns}
\end{frame}

% ---------------- PESSOA 2 ----------------
\section{Arquitetura MVT}

\begin{frame}{Arquitetura MVT}
  O Django adota o padrão \textbf{MVT} (\textit{Model -- View -- Template}):
  \vspace{0.3cm}

  \begin{columns}[T]
    \column{0.31\textwidth}
    \begin{block}{Model}
      Define a estrutura dos dados e comunica-se com o banco via ORM.\\[4pt]
      \footnotesize\textit{Ex.: classe \texttt{Tarefa}}
    \end{block}

    \column{0.31\textwidth}
    \begin{block}{View}
      Contém a regra de negócio: recebe a requisição HTTP e devolve a resposta.\\[4pt]
      \footnotesize\textit{Ex.: funções \texttt{listar}, \texttt{criar}}
    \end{block}

    \column{0.31\textwidth}
    \begin{block}{Template}
      Camada de apresentação em HTML dinâmico exibida ao usuário.\\[4pt]
      \footnotesize\textit{Ex.: \texttt{listar.html}}
    \end{block}
  \end{columns}
\end{frame}

\begin{frame}[fragile]{Estrutura de um Projeto Django}
  \begin{columns}[c]
    \column{0.45\textwidth}
\begin{lstlisting}
meu_projeto/
  manage.py
  meu_projeto/
    settings.py
    urls.py
  tarefas/
    models.py
    views.py
    urls.py
    templates/
\end{lstlisting}

    \column{0.50\textwidth}
    \begin{itemize}
      \setlength\itemsep{0.5em}
      \item \texttt{manage.py}: utilitário de linha de comando.
      \item \texttt{settings.py}: configurações globais e banco de dados.
      \item \texttt{urls.py}: mapeamento de rotas da aplicação.
      \item \texttt{tarefas/}: aplicativo modular com lógica própria.
    \end{itemize}
  \end{columns}
\end{frame}

% ---------------- PESSOA 3 ----------------
\section{Configuração}

\begin{frame}[fragile]{Instalação e Criação do Projeto}
\begin{lstlisting}
# 1. Criar e ativar o ambiente virtual
python -m venv venv
source venv/bin/activate

# 2. Instalar o Django
pip install django

# 3. Iniciar o projeto, criar o app e rodar o servidor
django-admin startproject meu_projeto
cd meu_projeto
python manage.py startapp tarefas
python manage.py runserver
\end{lstlisting}
\end{frame}

% ---------------- PESSOA 4 ----------------
\section{CRUD na Prática}

\begin{frame}{Nosso CRUD: Lista de Tarefas}
  Implementação das quatro operações básicas sobre o modelo \textbf{Tarefa}:
  \vspace{0.3cm}

  \begin{columns}[T]
    \column{0.48\textwidth}
    \begin{block}{\textbf{C}reate \& \textbf{R}ead}
      \begin{itemize}
        \item \textbf{Create:} cadastrar nova tarefa via formulário.
        \item \textbf{Read:} consultar e listar tarefas cadastradas.
      \end{itemize}
    \end{block}

    \column{0.48\textwidth}
    \begin{block}{\textbf{U}pdate \& \textbf{D}elete}
      \begin{itemize}
        \item \textbf{Update:} editar título ou marcar como concluída.
        \item \textbf{Delete:} remover tarefa do banco de dados.
      \end{itemize}
    \end{block}
  \end{columns}
\end{frame}

\begin{frame}[fragile]{1. Model --- \texttt{models.py}}
\begin{lstlisting}
from django.db import models

class Tarefa(models.Model):
    titulo = models.CharField(max_length=100)
    concluida = models.BooleanField(default=False)

    def __str__(self):
        return self.titulo
\end{lstlisting}
  \vspace{0.2cm}
  \small\color{textMuted} O ORM do Django converte essa classe Python diretamente em uma tabela no banco de dados.
\end{frame}

\begin{frame}[fragile]{2. View --- \texttt{views.py}}
\begin{lstlisting}
from django.shortcuts import render, redirect
from .models import Tarefa

def listar(request):
    tarefas = Tarefa.objects.all()
    return render(request, 'listar.html', {'tarefas': tarefas})

def criar(request):
    if request.method == 'POST':
        Tarefa.objects.create(titulo=request.POST['titulo'])
        return redirect('listar')
    return render(request, 'criar.html')
\end{lstlisting}
\end{frame}

\begin{frame}[fragile]{3. Rotas --- \texttt{urls.py}}
\begin{lstlisting}
from django.urls import path
from . import views

urlpatterns = [
    path('', views.listar, name='listar'),
    path('criar/', views.criar, name='criar'),
    path('editar/<int:id>/', views.editar, name='editar'),
    path('excluir/<int:id>/', views.excluir, name='excluir'),
]
\end{lstlisting}
\end{frame}

% ---------------- ENCERRAMENTO ----------------
\section{Encerramento}

\begin{frame}{Demonstração ao Vivo}
  \begin{center}
    \vspace{0.6cm}
    {\Huge\bfseries\color{djangoDark} Hora da Demonstração!}\\[0.5cm]
    {\large\color{textMuted} Execução no navegador do fluxo completo:}\\[0.4cm]
    {\large\bfseries\color{djangoMid} Criar \quad $\rightarrow$ \quad Listar \quad $\rightarrow$ \quad Editar \quad $\rightarrow$ \quad Excluir}
  \end{center}
\end{frame}

\begin{frame}{Conclusão}
  \begin{itemize}
    \setlength\itemsep{0.55em}
    \item O Django acelera o desenvolvimento web sem abrir mão da organização.
    \item A arquitetura \textbf{MVT} separa claramente dados, regras de negócio e interface.
    \item Com poucas linhas de código é possível construir um CRUD seguro e funcional.
    \item Excelente escolha tanto para aprendizado quanto para aplicações em produção.
  \end{itemize}
  \vspace{0.6cm}
  \begin{center}
    {\Large\bfseries\color{djangoDark} Obrigado! Perguntas?}
  \end{center}
\end{frame}

\end{document}